% Build (from week_6/module6, after the three commands in README.md):
%   pdflatex answers.tex && pdflatex answers.tex
\documentclass[10pt,letterpaper]{article}
\usepackage[margin=0.75in]{geometry}
\usepackage{amsmath,amssymb,graphicx,booktabs,float,xurl,hyperref}
\hypersetup{colorlinks=true,urlcolor=blue,linkcolor=black,citecolor=black}
\setlength{\parindent}{0pt}
\setlength{\parskip}{4pt}
\begin{document}

\begin{center}
{\Large CSc 8830 Assignment 6: Motion Tracking and Structure from Motion}\\[4pt]
Raghuveer Gummadi
\end{center}

\textbf{Live app:} \url{https://ragman95.pythonanywhere.com/} \quad (Module 6: \url{https://ragman95.pythonanywhere.com/module-6/})\\
\textbf{Code:} \url{https://github.com/RagDoll95/CSC8830_Assignments/tree/main/week_6/module6}

The app runs the repository scripts: \texttt{motion.py} (flow video), \texttt{validate.py}
(two-frame tracking check and flow statistics) and \texttt{sfm.py} (four-view factorization).
Theory and method follow the course lectures [N1] \emph{Optical Flow} (FPCV-4-3),
[N2] \emph{Structure from Motion} (FPCV-4-4) and [N3] \emph{Object Tracking} (FPCV-5-1);
slide numbers are given as, e.g., [N1, s.17].

%======================================================================
\section{Optical flow and motion tracking}

\textbf{Videos.} Penguin Walk, 00:30--01:00 [V1], and Sand Marble Race, 02:00--02:30 [V2]:
30\,s each, $640\times360$, 25 and 50 frames/s. In every frame up to 100 corners are
tracked into the next frame with coarse-to-fine Lucas--Kanade [N1, s.26--36] (OpenCV's
implementation of [Bouguet 2000], $21\times21$ window, 4 pyramid levels). The arrows form
the flow video; all vectors are in \texttt{points.csv}.

\subsection{What optical flow tells us, with evidence}
Optical flow is the motion of brightness patterns in the image; ideally it equals the
\emph{motion field}, the projection of 3-D scene motion [N1, s.3--6]. Each vector's length
is the point's image speed and its direction is the direction of motion [N1, s.6].
From our videos:
\begin{enumerate}\itemsep0pt
\item \textbf{Object motion: speed and direction.} The white marble moves
$(-2.0,-3.8)$\,px/frame ($\approx214$\,px/s, up-left) and the yellow one $(-0.4,+1.0)$\,px/frame.
Without the 3-D geometry of the track these stay in pixels; with a known ground plane
they could be converted to real speeds, as in traffic monitoring [N1, s.44].
\item \textbf{Camera motion (dominant flow).} The flow shared by most points is the
background, i.e.\ the camera's motion; removing it stabilizes video [N1, s.46]. Taking each
frame's median vector as that dominant flow: the hand-held penguin camera moves by a median
$3.8$\,px/frame (only 1\% of vectors are static), the marble camera pans smoothly at
$2.7$\,px/frame, about $(-4.3,1.7)$ in the first pair (plots below).
\item \textbf{Independently moving objects.} Vectors that disagree with the dominant flow
(magenta) mark things moving on their own: in the penguin pair they lie on the penguins and
the keeper's legs and hands, while the railing and boards follow the camera (cyan).
\item \textbf{Relative depth.} For a translating camera, nearer points move faster
[N1, s.6]. In the marble race all sand moves the same way, but the median speed falls from
5.3 (left third) to 3.9 (middle) to 2.9\,px/frame (right third) over the first 4\,s:
consistent with the sand being nearer on the left. (Camera rotation can also make flow
non-uniform, so this is relative, not metric, depth.)
\item \textbf{Where flow differs from the motion field, and where it fails.} The race's
scoreboard and timer overlays give zero flow (31\% of vectors are static) although the camera
moves: flow measures brightness motion, not scene motion [N1, s.7--8]. Spikes in the plots
are scene cuts (marbles 19.6\,s, penguins 24.1\,s) and a fast pan (penguins 7.2--7.5\,s,
over 30\,px/frame), where the small-motion assumption breaks [N1, s.34].
\end{enumerate}
\begin{figure}[H]\centering
\includegraphics[width=\linewidth]{output/penguins/flow_summary.png}\\
\includegraphics[width=\linewidth]{output/marble_race/flow_summary.png}
\caption{Left: frame 0$\to$1 flow (vectors $\times4$); cyan follows the frame's median
(dominant) flow, magenta differs by more than 1\,px. Right: dominant flow over the 30\,s and
the 90th percentile of deviation from it. Top: penguins; bottom: marbles.}
\end{figure}

\subsection{Tracking equations from first principles}
\textbf{The tracking problem for two frames.} Given the location of a target (a point and
a small window $W$ around it) in the frame at time $t$, find its location in the frame at
$t+\delta t$ [N3, s.22]. Template matching answers this by minimizing the sum of squared
differences over a shift $(\delta x,\delta y)$ [N3, s.25; N1, s.40]:
\[
\mathrm{SSD}(\delta x,\delta y)=\sum_{(x,y)\in W}\big[I(x+\delta x,\,y+\delta y,\,t+\delta t)-I(x,y,t)\big]^2 .
\]
This is slow over a large search and limited to whole pixels. The optical flow equations
give the minimizer directly, at sub-pixel accuracy.

\textbf{Assumption 1, brightness constancy} [N1, s.13]:
$I(x+\delta x,\,y+\delta y,\,t+\delta t)=I(x,y,t)$.

\textbf{Assumption 2, small motion} [N1, s.14--16]. The first-order Taylor expansion is
\[
I(x+\delta x,\,y+\delta y,\,t+\delta t)\approx I(x,y,t)+I_x\delta x+I_y\delta y+I_t\delta t .
\]
Subtracting the two, dividing by $\delta t$ and letting $\delta t\to0$ gives the
\textbf{optical flow constraint equation} [N1, s.17]
\[
I_xu+I_yv+I_t=0,\qquad (u,v)=\Big(\frac{dx}{dt},\frac{dy}{dt}\Big).
\]
$I_x,I_y,I_t$ come from two frames with the $2\times2\times2$ pixel cube [N1, s.18], e.g.
$I_x(i,j,t)=\tfrac14\big[I(i{+}1,j,t)+I(i{+}1,j,t{+}1)+I(i{+}1,j{+}1,t)+I(i{+}1,j{+}1,t{+}1)\big]
-\tfrac14\big[I(i,j,t)+I(i,j,t{+}1)+I(i,j{+}1,t)+I(i,j{+}1,t{+}1)\big]$.
One equation in two unknowns only fixes the \emph{normal flow},
$|u_n|=|I_t|/\sqrt{I_x^2+I_y^2}$ along $(I_x,I_y)$: the aperture problem [N1, s.19--24].

\textbf{Lucas--Kanade} [N1, s.26--29]. Assume $(u,v)$ is constant in the $n\times n$ window
$W$. Each pixel $(k,l)\in W$ gives $I_x(k,l)u+I_y(k,l)v=-I_t(k,l)$, so
\[
\underbrace{\begin{bmatrix}I_x(1,1)&I_y(1,1)\\\vdots&\vdots\\I_x(n,n)&I_y(n,n)\end{bmatrix}}_{A\ (n^2\times2)}
\underbrace{\begin{bmatrix}u\\v\end{bmatrix}}_{\mathbf u}
=\underbrace{\begin{bmatrix}-I_t(1,1)\\\vdots\\-I_t(n,n)\end{bmatrix}}_{B\ (n^2\times1)},
\qquad
\underbrace{\begin{bmatrix}\sum I_x^2&\sum I_xI_y\\\sum I_xI_y&\sum I_y^2\end{bmatrix}}_{A^TA}
\begin{bmatrix}u\\v\end{bmatrix}=
\underbrace{\begin{bmatrix}-\sum I_xI_t\\-\sum I_yI_t\end{bmatrix}}_{A^TB},
\]
$\mathbf u=(A^TA)^{-1}A^TB$. This is exactly the minimum of the linearized SSD above,
$\sum_W(I_xu+I_yv+I_t)^2$. It is reliable only when $A^TA$ is well conditioned: both
eigenvalues large and neither much larger than the other, i.e.\ a textured patch, not a
smooth region or a single edge [N1, s.29--32]. The tracked point is $p_1=p_0+\mathbf u$.

\textbf{Large motion} [N1, s.34--36]. When the motion is more than about a pixel, the
Taylor step fails. Build a pyramid by averaging $2\times2$ blocks (motion halves at each
level), solve at the coarsest level, then at each finer level warp the earlier frame by the
current flow, solve for the small remaining $(\Delta u,\Delta v)$ and add it.

\subsection{Bilinear interpolation}
The warp moves pixels to non-integer positions, so intensities must be sampled between
pixel centres. Let $x=m+\alpha$, $y=n+\beta$ with $m=\lfloor x\rfloor$, $n=\lfloor y\rfloor$,
$0\le\alpha,\beta<1$. Interpolate linearly along $x$ in the two neighbouring rows,
$I(x,n)=(1-\alpha)I_{m,n}+\alpha I_{m+1,n}$ and
$I(x,n{+}1)=(1-\alpha)I_{m,n+1}+\alpha I_{m+1,n+1}$, then along $y$:
\[
I(x,y)=(1-\beta)I(x,n)+\beta I(x,n{+}1)
=(1-\alpha)(1-\beta)I_{m,n}+\alpha(1-\beta)I_{m+1,n}+(1-\alpha)\beta I_{m,n+1}+\alpha\beta I_{m+1,n+1}.
\]
Each weight is the area of the opposite sub-rectangle; the four sum to 1, and the result
is exact for any function $a+bx+cy+dxy$. Example: neighbours
$\begin{bmatrix}10&20\\30&40\end{bmatrix}$, $\alpha=0.25$, $\beta=0.5$: rows $12.5$ and $32.5$,
so $I=22.5$.

\subsection{Validation against actual pixel locations}
For the first pair of consecutive frames in each video I marked recognizable points in
frame~0 and found them again in frame~1 by eye on enlarged images (integer pixels, about
$\pm1$\,px; \texttt{measurements.csv}). These are compared with the Lucas--Kanade prediction
and with SSD template matching ($21\times21$ template, $\pm10$\,px search).

\begin{table}[H]\centering\small
\begin{tabular}{llrrrrrr}\toprule
Video & Point & $p_0$ (frame 0) & Measured $p_1$ & LK $\hat p_1$ & LK error & SSD match & SSD error\\\midrule
Penguins & Neck patch   & (79, 49)  & (78, 50)  & (78.23, 49.96)  & 0.24 & (78, 50) & 0.00\\
Penguins & Beak tip     & (114, 14) & (115, 12) & (111.63, 13.85) & 3.84 & (111, 14) & 4.47\\
Marbles  & White marble & (288, 102)& (286, 99) & (285.98, 98.25) & 0.75 & (286, 98) & 1.00\\
Marbles  & Blue marble  & (432, 53) & (431, 51) & (430.79, 52.68) & 1.69 & (431, 52) & 1.00\\
Marbles  & Yellow marble& (428, 66) & (427, 67) & (427.60, 66.99) & 0.60 & (428, 67) & 1.00\\\bottomrule
\end{tabular}
\end{table}

\textbf{Worked solve (penguin neck).} With cube derivatives over the $21\times21=441$-pixel
window at $(79,49)$:
\[
\begin{bmatrix}130018.6&-20690.0\\-20690.0&47630.9\end{bmatrix}
\begin{bmatrix}u\\v\end{bmatrix}=\begin{bmatrix}-119434.5\\57979.9\end{bmatrix},
\quad \det(A^TA)=5.765\times10^{9},
\]
\[
\begin{bmatrix}u\\v\end{bmatrix}=\frac{1}{\det}
\begin{bmatrix}47630.9&20690.0\\20690.0&130018.6\end{bmatrix}
\begin{bmatrix}-119434.5\\57979.9\end{bmatrix}
=\begin{bmatrix}-0.779\\0.879\end{bmatrix}\text{ px/frame},
\]
so one step predicts $(78.22,49.88)$, $0.25$\,px from the measured $(78,50)$. The
eigenvalues of $A^TA$ are $4.3\times10^4$ and $1.3\times10^5$: a well-conditioned,
textured patch. At the beak tip they are only $2.3\times10^3$ and $8.1\times10^3$ (blurred,
low contrast), and both LK and SSD miss by about 4\,px. For the white marble one linear
step gives $(-0.68,-1.24)$ and misses by $2.2$\,px, because the $\approx4$\,px motion breaks
the small-motion assumption; coarse-to-fine iteration reduces this to $0.75$\,px. All five
systems are in \texttt{tracking\_calculations.json}.

\begin{figure}[H]\centering
\includegraphics[width=.49\linewidth]{output/penguins/validation.png}\hfill
\includegraphics[width=.49\linewidth]{output/marble_race/validation.png}
\caption{Measured points (green circles) and Lucas--Kanade predictions (red $+$) in
frames 0 and 1. Left: penguins; right: marbles.}
\end{figure}

%======================================================================
\section{Structure from motion with four views of a planar object}

\textbf{Object, images and camera.} The object is a flat $9\times6$-corner checkerboard
(24.9\,mm squares) shown on a laptop screen, photographed from four directions with an
iPhone (1$\times$ main lens, $3024\times4032$ portrait). Its 54 corners are detected
automatically and are the tracked features [N2, s.5]. The camera parameters come from my
Module 2 calibration (RMS reprojection 1.86\,px):
\[
K=\begin{bmatrix}3336.69&0&1496.68\\0&3335.58&2020.27\\0&0&1\end{bmatrix},\quad
(k_1,k_2,p_1,p_2,k_3)=(0.223,-1.376,0.0006,-0.0008,2.709).
\]
They are used only to remove lens distortion and convert pixels to normalized image
coordinates $K^{-1}[u,v,1]^T$. The camera positions below come from the same calibration
(the board's pose in each photo). The factorization does \emph{not} use them; they are used
only to check it. Tomasi--Kanade assumes one constant magnification [N2, s.8], so the
four photos were chosen at nearly the same distance (1035--1072\,mm) and from different
directions.

\begin{table}[H]\centering\small
\begin{tabular}{clrrrr}\toprule
View & Photo & Camera position (mm, board frame) & Distance & Tilt from normal & Depth range on board\\\midrule
1 & IMG\_3930 & (44, $-389$, $-930$)   & 1035 mm & 26.0$^\circ$ & 7.9\%\\
2 & IMG\_3917 & ($-419$, $-125$, $-915$) & 1069 mm & 31.1$^\circ$ & 11.5\%\\
3 & IMG\_3928 & (631, $-502$, $-709$)  & 1050 mm & 47.5$^\circ$ & 16.2\%\\
4 & IMG\_3938 & ($-374$, $-538$, $-751$) & 1072 mm & 45.5$^\circ$ & 13.9\%\\\bottomrule
\end{tabular}
\end{table}
\begin{figure}[H]\centering
\includegraphics[width=\linewidth]{output/structure/views.png}
\caption{The four views. Green: detected corners; red: reprojection of the factorization
$MS$; yellow: estimated boundary.}
\end{figure}

\subsection{Method (Tomasi--Kanade factorization [N2])}
\textbf{Orthographic projection} [N2, s.8--11]. Frame $f$ has image axes $\mathbf i_f,\mathbf j_f$
(unit, orthogonal) and position $C_f$; scene point $P_p$ projects to
$u_{f,p}=\mathbf i_f^T(P_p-C_f)$, $v_{f,p}=\mathbf j_f^T(P_p-C_f)$.

\textbf{Centering trick} [N2, s.12--14]. Put the world origin at the centroid of the points.
Then the image centroid is $\bar u_f=-\mathbf i_f^TC_f$, and the centroid-subtracted coordinates
$\tilde u_{f,p}=u_{f,p}-\bar u_f=\mathbf i_f^TP_p$, $\tilde v_{f,p}=\mathbf j_f^TP_p$ no longer
contain $C_f$. Stacking all $F=4$ frames and $N=54$ points:
\[
\underbrace{\begin{bmatrix}\tilde u_{1,1}&\cdots&\tilde u_{1,N}\\\vdots&&\vdots\\\tilde u_{F,1}&\cdots&\tilde u_{F,N}\\
\tilde v_{1,1}&\cdots&\tilde v_{1,N}\\\vdots&&\vdots\\\tilde v_{F,1}&\cdots&\tilde v_{F,N}\end{bmatrix}}_{W\ (2F\times N)\ \text{known}}
=\underbrace{\begin{bmatrix}\mathbf i_1^T\\\vdots\\\mathbf i_F^T\\\mathbf j_1^T\\\vdots\\\mathbf j_F^T\end{bmatrix}}_{M\ (2F\times3)}
\underbrace{\begin{bmatrix}P_1&\cdots&P_N\end{bmatrix}}_{S\ (3\times N)} .
\]

\textbf{Rank} [N2, s.24--27]. $\mathrm{rank}(W)=\mathrm{rank}(MS)\le\min(\mathrm{rank}\,M,\mathrm{rank}\,S)\le3$.
For a \emph{planar} object, take the world $z$ axis along the plane normal, so every
$P_p=(x_p,y_p,0)$ and $\mathrm{rank}(S)=2$. Then $\mathrm{rank}(W)\le2$, and only the in-plane
parts $\tilde{\mathbf i}_f,\tilde{\mathbf j}_f$ (first two components) of the camera axes appear in $W$:
$W=\tilde M\tilde S$ with $\tilde M$ $2F\times2$ and $\tilde S$ $2\times N$.

\textbf{SVD and factorization} [N2, s.28--33]. $W=U\Sigma V^T$. Keep the two largest singular
values: $W=U_1\Sigma_1V_1^T$, and set $\tilde M=U_1\Sigma_1^{1/2}Q$,
$\tilde S=Q^{-1}\Sigma_1^{1/2}V_1^T$ for an unknown $2\times2$ matrix $Q$.

\textbf{Orthonormality} [N2, s.34--35]. Write $\mathbf i_f=(\hat a_f^TQ,\ \alpha_f)$ and
$\mathbf j_f=(\hat b_f^TQ,\ \beta_f)$, where $\hat a_f^T,\hat b_f^T$ are rows $f$ and $F+f$ of
$U_1\Sigma_1^{1/2}$. The out-of-plane components $\alpha_f,\beta_f$ multiply $z=0$, so they are
invisible in $W$, but $\mathbf i_f,\mathbf j_f$ must still be orthonormal:
\[
\hat a_f^TQQ^T\hat a_f+\alpha_f^2=1,\qquad \hat b_f^TQQ^T\hat b_f+\beta_f^2=1,\qquad
\hat a_f^TQQ^T\hat b_f+\alpha_f\beta_f=0,\qquad f=1..F.
\]
These are $3F=12$ quadratic equations in the 4 entries of $Q$ and the 8 values $\alpha_f,\beta_f$.
As on [N2, s.35], they are solved with Newton's method (Gauss--Newton, 40 deterministic
starts, best least-squares solution kept). $Q$ is defined only up to an in-plane rotation and
a mirror image. I fix the rotation by putting the first edge on $+x$, and the mirror by
requiring that every camera sees the front of the plane, $(\mathbf i_f\times\mathbf j_f)_z>0$.

\textbf{Scale and boundary.} Orthography gives shape up to scale. With normalized coordinates
$u\approx X/\bar Z$, $\tilde S$ is in units of $1/\bar Z$. One known length fixes it: the top
edge (corner 1 to 9) is $8\times24.9=199.2$\,mm. The boundary is the polygon through the
outer corners of $\tilde S$.

\subsection{Calculations}
$W$ is $8\times54$. Its columns for corners 1, 9, 46 and 54 (the four outer corners) are
\[
W_{[:,\,1,9,46,54]}=\begin{bmatrix}
-0.1011&0.1062&-0.0973&0.0926\\-0.0893&0.0792&-0.0792&0.0824\\-0.0756&0.0818&-0.0715&0.0750\\-0.0834&0.0836&-0.0808&0.0738\\
-0.0592&-0.0585&0.0555&0.0567\\-0.0472&-0.0722&0.0747&0.0433\\-0.0794&-0.0208&0.0175&0.0828\\-0.0247&-0.0730&0.0735&0.0223
\end{bmatrix},
\]
(rows $\tilde u_1..\tilde u_4$, $\tilde v_1..\tilde v_4$; image centroids e.g.\ $(\bar u_1,\bar v_1)=(0.0545,-0.2421)$).
\[
\sigma(W)=(0.8304,\ 0.5425,\ 0.0159,\ 0.0131,\ 0.0029,\ 0.0005,\dots).
\]
Two singular values dominate, about $35\times$ larger than the rest, confirming rank 2 for
a plane; the small remainder is perspective and corner noise. Then
\[
U_1\Sigma_1^{1/2}=\begin{bmatrix}
-0.5165&-0.0072\\-0.4282&-0.0422\\-0.3947&-0.0150\\-0.4182&-0.0061\\
0.0092&-0.3915\\0.0856&-0.4014\\-0.1512&-0.3484\\0.1391&-0.3235\end{bmatrix},\qquad
Q=\begin{bmatrix}-2.0015&0.0446\\-0.0202&-2.4469\end{bmatrix},\qquad
QQ^T=\begin{bmatrix}4.0082&-0.0688\\-0.0688&5.9875\end{bmatrix}.
\]
The recovered camera axes $\mathbf i_f=(\tilde{\mathbf i}_f,\alpha_f)$, $\mathbf j_f=(\tilde{\mathbf j}_f,\beta_f)$,
$\mathbf k_f=\mathbf i_f\times\mathbf j_f$, and tilt $=\arccos(k_{f,z}/\|\mathbf k_f\|)$:
\begin{table}[H]\centering\small
\begin{tabular}{crrrrr}\toprule
View & $\mathbf i_f$ & $\mathbf j_f$ & $\mathbf k_f$ & Recovered tilt & Calibration tilt\\\midrule
1 & (1.034, $-0.005$, 0.054) & ($-0.011$, 0.958, $-0.300$) & ($-0.050$, 0.310, 0.991) & 17.6$^\circ$ & 26.0$^\circ$\\
2 & (0.858, 0.084, 0.508)   & ($-0.163$, 0.986, 0.104)  & ($-0.492$, $-0.172$, 0.860) & 31.2$^\circ$ & 31.1$^\circ$\\
3 & (0.790, 0.019, 0.612)   & (0.310, 0.846, $-0.432$)  & ($-0.525$, 0.531, 0.662) & 48.4$^\circ$ & 47.5$^\circ$\\
4 & (0.837, $-0.004$, 0.529) & ($-0.272$, 0.798, 0.517) & ($-0.424$, $-0.576$, 0.667) & 47.0$^\circ$ & 45.5$^\circ$\\\bottomrule
\end{tabular}
\end{table}
Check for view 2: $|\mathbf i_2|^2=0.858^2+0.084^2+0.508^2=1.001$ and
$\mathbf i_2\cdot\mathbf j_2=-0.140+0.083+0.053=-0.004\approx0$. The largest equation residual is
$0.07$ (view 1, $|\mathbf i_1|=1.034$). View 1, the least tilted, has the smallest out-of-plane
components, so its tilt is the least well determined (off by 8$^\circ$). For a plane only
$\alpha_f\beta_f$ is fixed, so each camera's tilt angle is recovered but not which way it
leans. The camera positions $C_f$ were removed by the centering trick, so orthographic
factorization recovers orientations, not positions [N2, s.14].

\textbf{Structure.} Corner 1 is $\tilde S_1=(-0.0972,-0.0624)$ and corner 9 is
$(0.1008,-0.0588)$, $0.1980$ apart, so the scale is $199.2/0.1980=1006$\,mm per unit.
Because normalized coordinates are $\approx X/\bar Z$, this scale should equal the mean camera
distance, and it does: 1006\,mm vs.\ 1057\,mm from calibration. Corner 1 is
then $(-97.7,-62.8)$\,mm. Reprojection check, corner 1 in view 2:
$\tilde{\mathbf i}_2^T\tilde S_1=0.8579(-0.0972)+0.0841(-0.0624)=-0.0886$ vs.\ the measured
$\tilde u_{2,1}=-0.0893$: $0.0007\times3337\approx2.3$\,px.

\begin{figure}[H]\centering
\includegraphics[width=.95\linewidth]{output/structure/reconstruction.png}
\caption{Left: singular values of $W$. Right: recovered structure and boundary (blue),
with the true grid rigidly aligned (grey dashed).}
\end{figure}

\textbf{Result.} The recovered boundary sides are 199.2 (scale reference), 120.5, 187.4 and
121.2\,mm, against the true 199.2 and 124.5\,mm. All 54 points are within 7.7\,mm of the true
grid after rigid alignment (RMS 2.9\,mm). Mean reprojection error is 2.0--6.4\,px per view.
Three of four camera tilts agree with the calibration to within 1.5$^\circ$.

\textbf{Why the errors are not zero.} The orthographic model needs the depth variation to be
small compared with the distance [N2, s.8]. Here the tilted board spans 8--16\% of the
viewing distance, so perspective foreshortens the far edge differently in each view. The
model cannot represent that, which is why the bottom edge comes out 187\,mm instead of
199\,mm. The four photos are also among those used for the Module 2 calibration, so $K$ is
fitted to them. Larger distances (with zoom) or the perspective extensions mentioned in
[N2, s.38] would reduce the error. The app accepts any four new photos with corners clicked
by hand.

%======================================================================
\section*{References}
\small
\begin{itemize}\itemsep0pt
\item[{[N1]}] S.~K. Nayar, \emph{Optical Flow}, Monograph FPCV-4-3, First Principles of Computer Vision, Columbia University, April 2025.
\item[{[N2]}] S.~K. Nayar, \emph{Structure from Motion}, Monograph FPCV-4-4, First Principles of Computer Vision, Columbia University, April 2025.
\item[{[N3]}] S.~K. Nayar, \emph{Object Tracking}, Monograph FPCV-5-1, First Principles of Computer Vision, Columbia University, April 2025.
\item[] B.~D. Lucas and T.~Kanade, ``An iterative image registration technique with an application to stereo vision,'' Imaging Understanding Workshop, 1981.
\item[] B.~K.~P. Horn and B.~G. Schunck, ``Determining optical flow,'' \emph{Artificial Intelligence}, 1981.
\item[] J.-Y. Bouguet, ``Pyramidal implementation of the Lucas Kanade feature tracker,'' Intel Corp., 2000 (used via OpenCV).
\item[] C.~Tomasi and T.~Kanade, ``Shape and motion from image streams under orthography: a factorization method,'' \emph{IJCV}, 1992.
\item[] Z.~Zhang, ``A flexible new technique for camera calibration,'' \emph{IEEE TPAMI}, 22(11), 2000 (Module 2 calibration).
\item[{[V1]}] ``Penguin Walk -- Zoo Basel [HD],'' \url{https://www.youtube.com/watch?v=fdk7X4LrgS4}, sample 00:30--01:00.
\item[{[V2]}] ``SAND MARBLE RACE, S7: Round 1,'' Jelle's Marble Runs, \url{https://www.youtube.com/watch?v=0sHu99vdz-A}, sample 02:00--02:30.
\end{itemize}
\end{document}
